%% explanation_v3_improved.tex — technical specification, current implementation.
% !TEX program = pdflatex

\documentclass[11pt, a4paper]{article}
\usepackage{amsmath, amssymb}
\usepackage{booktabs}
\usepackage{geometry}
\usepackage{microtype}
\usepackage{parskip}
\usepackage[hidelinks]{hyperref}
\usepackage{algorithm}
\usepackage{algpseudocode}
\geometry{margin=2.5cm}

\title{\textbf{v3\_improved Environment: Technical Specification}}
\author{}
\date{}

\begin{document}
\maketitle

% ─────────────────────────────────────────────────────────────────────────────
\section{Overview}
% ─────────────────────────────────────────────────────────────────────────────

The environment simulates a single upper-airspace sector populated with A320
aircraft cruising at FL350 ($M = 0.78$, approximately 450\,kt TAS).
BlueSky is used as the underlying physics engine: it integrates aircraft
kinematics at a fixed timestep of 0.5\,s, handling bank-angle turns,
Mach-number commands, and sector-boundary detection.
The reinforcement learning agent acts as an air traffic controller and issues
one clearance per step to the focus aircraft selected by the urgency matrix.

\paragraph{Episode initialisation.}
At the start of each episode, the number of aircraft is sampled uniformly as
$N \sim \mathcal{U}\{2,\ldots,15\}$ and the sector density as
$\rho \sim \mathcal{U}(5\,000,\,15\,000)$\,km$^2$/aircraft.
The sector area is set to $N \times \rho$, which ensures a uniform joint
distribution over both $N$ and $\rho$ without conflating the two.
A random convex polygon with 5--7 vertices is generated and scaled to the
target area; polygons with an isoperimetric ratio below 0.65 are rejected to
avoid degenerate elongated shapes.
Each aircraft is spawned on a distinct arc of the sector perimeter with
uniform position jitter, flying toward a reference point on the opposite arc
and then toward a destination two sector diameters beyond the boundary.
Minimum spawn separation is 10\,NM.
Aircraft that exit the sector are replaced within one step to maintain constant
traffic density throughout the episode.
Episode length is approximately 60\,min of simulated time (2.5 sector crossings),
computed from the sector diameter and cruise speed.

\paragraph{Step granularity.}
Each RL step spans $10 \times 0.5\,\text{s} = 5\,\text{s}$ of simulated time.
At the standard A320 bank angle a $45^\circ$ heading change takes approximately
40\,s (8 steps), so the geometric consequence of a clearance is observable
within two to three steps.
Speed commands settle over 2--5 steps depending on the Mach delta and
deceleration rate.

% ─────────────────────────────────────────────────────────────────────────────
\section{Environment Loop}
\label{sec:loop}
% ─────────────────────────────────────────────────────────────────────────────

\begin{algorithm}[h]
\caption{\textsc{Reset} — initialise episode}
\label{alg:reset}
\begin{algorithmic}[1]
\State $N \sim \mathcal{U}\{2,\ldots,15\}$;\quad $\rho \sim \mathcal{U}(5\,000,\,15\,000)$;\quad area $\leftarrow N\times\rho$
\State Generate convex polygon; scale to area; reject if circularity $< 0.65$
\State Compute $T_\text{max}$ from sector diameter and cruise speed
\For{slot $s = 1,\ldots,N$} spawn aircraft on perimeter arc $s$ with jitter \EndFor
\State Issue \texttt{ASAS OFF}; define sector area filter
\State Compute urgency matrix $\mathbf{U}$
\State $f \leftarrow \arg\max_i \max_{j\neq i}\, u_{ij}$ \quad (highest-urgency aircraft)
\State \Return observation $\mathbf{o}(f)$
\end{algorithmic}
\end{algorithm}

\begin{algorithm}[h]
\caption{\textsc{Step}(action) — one RL step, 5\,s simulated}
\label{alg:step}
\begin{algorithmic}[1]
\State Spawn any aircraft whose replacement delay has elapsed
\State \textbf{Apply action to focus aircraft $f$ immediately:}
\If{heading offset $\delta$}
    \State $\psi_\text{cmd} \leftarrow \psi_\text{dest}(f) + \delta$ \quad\textit{(bearing recomputed from current position)}
    \If{$\delta = 0$ and $|\psi_\text{cmd} - \psi_\text{actual}| > 5^\circ$}
        \State \textbf{skip} \quad\textit{(mid-turn guard: direct action dropped)}
    \EndIf
\ElsIf{Mach adjustment $\Delta M$}
    \State $M_\text{cmd} \leftarrow \operatorname{clip}(M_\text{cmd} + \Delta M,\;0.70,\;0.82)$
\ElsIf{hold ($\varnothing$)}
    \State Re-issue $\psi_\text{cmd}$ unchanged \quad\textit{(reinforces current commanded heading)}
\EndIf
\State \textbf{Run physics:} advance BlueSky $10 \times 0.5\,\text{s}$ sub-steps uninterrupted
\State At sub-step 10: check LoS — $\ell_\text{LoS} \leftarrow \mathbf{1}[\exists\,j : d(f,j) < d_\text{sep}]$
\State Process exits; schedule replacements
\State Recompute $\mathbf{U}$; select new focus $f'$ (Section~\ref{sec:urgency})
\If{$f' \neq f$} mark grace step (Section~\ref{sec:reward}) \EndIf
\State $r \leftarrow r_\text{conflict}(f) + r_\text{LoS}(f) + r_\text{nav}(f) - w_\text{work}\cdot c_f$
\State \Return $\mathbf{o}(f'),\; r,\; \text{truncated} = (\text{step} \geq T_\text{max})$
\end{algorithmic}
\end{algorithm}

\paragraph{Action timing.}
The clearance is applied at the very start of the step, before any physics
sub-step runs.
The 10 sub-steps that follow execute uninterrupted: no further actions can be
issued during that window.
Loss-of-separation is evaluated only at the final sub-step, capturing whether
any separation violation occurred during the 5\,s window rather than merely
at a single instant.

\paragraph{Heading commands.}
The destination bearing $\psi_\text{dest}$ is recomputed from the aircraft's
current position at each step, so the ``direct'' action ($\delta = 0$) always
means fly to destination from where you are now, not where you were when the
conflict began.
A mid-turn guard prevents the direct action from being issued while the
aircraft is already executing a turn: if the difference between the commanded
and actual heading exceeds $5^\circ$, the direct action is silently dropped and
the previous commanded heading is retained.
This avoids issuing conflicting commands that would reset a partially completed
turn.

\paragraph{Observation at step boundary.}
The observation returned at the end of a step is constructed from the
\emph{new} focus aircraft $f'$, which may differ from $f$, the aircraft that
was controlled during that step.
The reward, however, is always attributed to $f$, the aircraft for which the
action was taken.
This separation is important: the policy conditions on $\mathbf{o}(f')$ when
selecting the next action, but the credit for the current step is assigned
entirely to $f$.
A grace step (Section~\ref{sec:reward}) handles the transition to ensure the
reward signal is well-defined immediately after a focus switch.

% ─────────────────────────────────────────────────────────────────────────────
\section{Urgency Matrix and Focus Selection}
\label{sec:urgency}
% ─────────────────────────────────────────────────────────────────────────────

\begin{equation}
    u_{ij} =
    \begin{cases}
        1 + 9\!\left(1 - \dfrac{d_{ij}}{d_\text{sep}}\right)
            & d_{ij} < d_\text{sep}
              \quad\text{(LoS,\;} u_{ij}\in[1,10]\text{)} \\[10pt]
        \dfrac{T_\text{warn} - t_\text{CPA}^{ij}}{T_\text{warn}}
            & d_\text{CPA}^{ij} < d_\text{sep},\;
              0 \leq t_\text{CPA}^{ij} \leq T_\text{look} \\[8pt]
        0   & \text{otherwise}
    \end{cases}
    \label{eq:urgency}
\end{equation}

$d_\text{sep}=5$\,NM,\quad $T_\text{warn}=600$\,s,\quad $T_\text{look}=900$\,s.

Focus: $f = \arg\max_i \max_{j\neq i} u_{ij}$.
No-conflict fallback: aircraft with largest heading deviation $(1-\cos\Delta\psi)/2$,
subject to a switch margin of 0.05 to prevent oscillation.

\paragraph{Commitment rule.}
Focus locked while $u_f>0$ or cooldown $< n_\text{clear}=5$ steps.
Emergency override: another aircraft reaching $u_\text{emerg}=0.8$
($t_\text{CPA}\approx120$\,s) releases the lock immediately.

% ─────────────────────────────────────────────────────────────────────────────
\section{Action Space}
% ─────────────────────────────────────────────────────────────────────────────

10 discrete actions (6 heading, 4 speed):
\begin{align}
    \psi_\text{cmd} &= \psi_\text{dest} + \delta,\qquad
    \delta \in \{-45^\circ,\,-30^\circ,\,0^\circ,\,+30^\circ,\,+45^\circ,\,\varnothing\} \\
    M_\text{cmd} &\leftarrow \operatorname{clip}(M_\text{cmd}+\Delta M,\,0.70,\,0.82),\qquad
    \Delta M \in \{-0.04,\,-0.02,\,+0.02,\,+0.04\}
\end{align}
Direct ($0^\circ$) and hold ($\varnothing$) are free ($c_f=0$); turns $c_f=1.0$; speed $c_f=0.5$.

% ─────────────────────────────────────────────────────────────────────────────
\section{Observation Space ($\mathbb{R}^{25}$)}
% ─────────────────────────────────────────────────────────────────────────────

Ego-centric from focus aircraft. Intruder slots sorted by urgency $\downarrow$
then CPA distance $\uparrow$; diverging pairs sentinel-sorted to back.
Empty/diverging slot: $(1,1,0,1,1)$.

\begin{table}[h]
\centering
\caption{Observation vector.}
\label{tab:obs}
\smallskip
\begin{tabular}{clll}
\toprule
\textbf{Index} & \textbf{Feature} & \textbf{Definition} & \textbf{Range} \\
\midrule
0 & $\sin\Delta\psi$ & $\sin(\psi_\text{dest}-\psi_\text{cmd})$ & $[-1,1]$ \\
1 & $\cos\Delta\psi$ & $\cos(\psi_\text{dest}-\psi_\text{cmd})$ & $[-1,1]$ \\
2 & Turn progress & $\operatorname{clip}((\psi_\text{cmd}-\psi_\text{act})/45^\circ,-1,1)$ & $[-1,1]$ \\
3 & Time to exit & $\min(d_\text{boundary}/(v\,T_\text{look}),1)$ & $[0,1]$ \\
4 & Speed offset & $(M_\text{cmd}-0.70)/0.12$ & $[0,1]$ \\[4pt]
\multicolumn{4}{l}{\textit{4 intruders $\times$ 5 features, sorted urgency $\downarrow$ / CPA $\uparrow$:}} \\
$+0$ & $\Delta\tilde{x}$ & lateral displacement / $d_\text{warn}$\,(75\,NM) & $[-1,1]$ \\
$+1$ & $\Delta\tilde{y}$ & longitudinal displacement / $d_\text{warn}$ & $[-1,1]$ \\
$+2$ & $\Delta\tilde{x}^\text{CPA}$ & lateral CPA offset / $d_\text{sep}$ & $[-1,1]$ \\
$+3$ & $\Delta\tilde{y}^\text{CPA}$ & longitudinal CPA offset / $d_\text{sep}$ & $[-1,1]$ \\
$+4$ & $\hat{t}^\text{CPA}$ & $\operatorname{clip}(t_\text{CPA}/T_\text{look},0,1)$ & $[0,1]$ \\
\bottomrule
\end{tabular}
\end{table}

Intruder positions in ownship heading frame, normalised by
$d_\text{warn} = T_\text{warn}\cdot v_\text{nom} = 75$\,NM (10-min action horizon).
CPA offset:
\begin{equation}
    \begin{pmatrix}\Delta x^\text{CPA}\\\Delta y^\text{CPA}\end{pmatrix}
        = \begin{pmatrix}\Delta x\\\Delta y\end{pmatrix}
        + t_\text{CPA}\begin{pmatrix}\Delta v_e\\\Delta v_n\end{pmatrix}.
\end{equation}

% ─────────────────────────────────────────────────────────────────────────────
\section{Reward Function}
\label{sec:reward}
% ─────────────────────────────────────────────────────────────────────────────

The total reward at each step is the sum of four components:
\begin{equation}
    r_t = r_\text{conflict} + r_\text{LoS} + r_\text{nav} - w_\text{work}\cdot c_f.
    \label{eq:reward_total}
\end{equation}
Each component addresses a distinct aspect of the controller's task.
The design follows an event-based rather than a continuous-proximity
philosophy: the conflict term fires on \emph{transitions} between safe and
unsafe states rather than penalising absolute separation distance at every
step.
This choice is motivated by the credit-assignment problem that arises when the
focus aircraft changes: if the reward were based on a continuous CPA distance,
a focus switch would cause the signal to jump to a completely different value
with no connection to the action just taken.
By anchoring the reward to discrete state transitions, all reward signals for
a given conflict sequence are guaranteed to be attributed to the aircraft that
was controlled throughout that sequence.

\paragraph{Conflict transition term $r_\text{conflict}$.}
Let $\sigma_t = \mathbf{1}[\max_{j \neq f} u_{fj} > 0]$ denote whether the
focus aircraft is in a predicted conflict at step $t$.
The conflict term is defined by the transition from step $t-1$ to step $t$:
\begin{equation}
    r_\text{conflict} =
    \begin{cases}
        +w_\text{resolve} & \sigma_{t-1}=1,\;\sigma_t=0 \quad\text{(conflict resolved)}\\
        -w_\text{resolve} & \sigma_{t-1}=0,\;\sigma_t=1 \quad\text{(new conflict)}\\
        -w_\text{resolve} & \sigma_{t-1}=1,\;\sigma_t=1 \quad\text{(ongoing conflict)}\\
        0 & \sigma_{t-1}=0,\;\sigma_t=0 \quad\text{(clean flight)}
    \end{cases}
    \label{eq:r_conflict}
\end{equation}
Three properties of this formulation are worth noting.
First, the reward is symmetric: the positive signal for resolving a conflict
($+w_\text{resolve}$) has the same magnitude as the penalty for allowing one
to develop ($-w_\text{resolve}$).
This prevents the agent from treating conflict avoidance as the only objective
and ignoring resolution once a conflict is already active.
Second, an ongoing conflict incurs a penalty at every step, not only when it
first appears.
This creates a sustained incentive to resolve conflicts quickly rather than
manoeuvring to the edge of the predicted separation zone and waiting.
Third, a clean step (no conflict before or after) carries zero reward, so
the conflict term does not penalise the agent simply for being in a busy
sector.

\paragraph{Grace step after focus switch.}
When the focus switches from aircraft $f$ to a new aircraft $f'$, the value
of $\sigma_{t-1}$ for $f'$ is not defined within the current episode context,
because the previous step was spent observing $f$.
On the first step of a new focus, $r_\text{conflict}$ is therefore set to zero
and $\sigma_{t-1}$ is initialised to $f'$'s actual current conflict state.
From the following step onward, the transition signal is computed normally.
Without this grace step, every focus switch would spuriously fire either a
``new conflict'' or a ``resolved'' signal, injecting noise into the gradient.

\paragraph{Loss-of-separation term $r_\text{LoS}$.}
A loss of separation is recorded if, at the final BlueSky sub-step, any
other aircraft is within $d_\text{sep} = 5$\,NM of the focus aircraft:
\begin{equation}
    r_\text{LoS} = -w_\text{LoS} \cdot \mathbf{1}[\ell_\text{LoS}].
    \label{eq:r_los}
\end{equation}
This term is additive: a step in which an LoS occurs receives both the ongoing
conflict penalty ($-w_\text{resolve}$) and the LoS penalty ($-w_\text{LoS}$),
giving a total of $-(w_\text{resolve} + w_\text{LoS}) = -4$ at that step.
The larger weight ($w_\text{LoS} = 3$) reflects the severity of an actual
separation violation compared with a merely predicted one.
Because LoS is checked only at the end of the 5\,s window, a very brief
incursion that begins and ends within the same step will still be detected
as long as it persists until sub-step 10.

\paragraph{Navigation term $r_\text{nav}$.}
\begin{equation}
    r_\text{nav} = -w_\text{nav} \cdot \frac{1 - \cos\Delta\psi_\text{act}}{2},
    \label{eq:r_nav}
\end{equation}
where $\Delta\psi_\text{act} = \psi_\text{dest} - \psi_\text{actual}$ is the
angular error between the aircraft's current actual heading and its destination
bearing.
This term is evaluated only when no conflict is active, that is, when
$\sigma_t = 0$.
It is suppressed during active conflict steps so that the navigation penalty
does not compete with the conflict signal during a resolution manoeuvre.
Its purpose is to prevent the agent from resolving a conflict by deviating onto
a heading that is never corrected: once the conflict is cleared, the agent
should return the aircraft to its planned route.
The weight $w_\text{nav} = 0.50$ is intentionally smaller than $w_\text{resolve}$
so that route adherence remains a secondary objective relative to separation.
The $(1 - \cos\Delta\psi)/2$ form maps the heading error to $[0, 1]$:
it is zero on direct track, $0.5$ when perpendicular, and $1$ when flying
directly away from the destination.

\paragraph{Workload term.}
The workload cost $c_f \in \{0, 0.5, 1.0\}$ reflects the operational cost of
the clearance issued.
Heading turns ($c_f = 1.0$) are more disruptive to route planning than speed
adjustments ($c_f = 0.5$); the direct and hold actions are free ($c_f = 0$).
The weight $w_\text{work} = 0.10$ is small enough that it never overrides
safety-driven actions, but it discourages issuing unnecessary instructions
when the traffic picture is already clear.

\begin{table}[h]
\centering
\caption{Reward weights.}
\label{tab:weights}
\smallskip
\begin{tabular}{llr}
\toprule
$w_\text{resolve}$ & Conflict transition (and ongoing conflict) & $1.00$ \\
$w_\text{LoS}$     & Loss-of-separation (additive)              & $3.00$ \\
$w_\text{nav}$     & Heading drift from destination (conflict-free steps only) & $0.50$ \\
$w_\text{work}$    & Workload per instruction                    & $0.10$ \\
\bottomrule
\end{tabular}
\end{table}

% ─────────────────────────────────────────────────────────────────────────────
\section{Training}
% ─────────────────────────────────────────────────────────────────────────────

PPO (Stable-Baselines3), 4 parallel environments, two-layer MLP.
Observations not normalised (\texttt{norm\_obs=False}): all features are
bounded by design.

\begin{table}[h]
\centering
\caption{PPO hyperparameters.}
\label{tab:hp}
\smallskip
\begin{tabular}{llr}
\toprule
$\gamma$               & Discount factor          & $0.99$ \\
$\lambda$              & GAE parameter            & $0.95$ \\
$T$                    & Steps per env per update & $1024$ \\
$N_\text{envs}$        & Parallel environments    & $4$    \\
$B$                    & Mini-batch size          & $64$   \\
$K$                    & Epochs per update        & $10$   \\
$\epsilon_\text{clip}$ & PPO clip range           & $0.2$  \\
$\beta_H$              & Entropy coefficient      & $0.02$ \\
$\alpha_\text{LR}$     & Learning rate            & $3\times10^{-4}$ \\
\texttt{net\_arch}     & MLP hidden layers        & $[128,\,128]$ \\
\bottomrule
\end{tabular}
\end{table}

\end{document}
